\section{The stationary edge input}\label{sec:background}

We fix the finite normalization and record the stationary edge,
static-comparison, and relaxed-energy inputs used below. Source cap
estimates are stated at their points of use in
Section~\ref{sec:edge-estimates}; the following sections develop the
extensions, entrance construction, and disorder comparison.

\subsection{The finite form and spectral coordinates}

Write \(\mu=\pi_{n,J}\). For a function \(f\) on the cube let
\[
 D_i f(\sigma)=\frac{f(\sigma^{i,+})-f(\sigma^{i,-})}{2},
 \qquad
 t_i(\sigma)=\tanh h_i(\sigma),\qquad
 v_i(\sigma)=1-t_i(\sigma)^2,
\]
where \(\sigma^{i,\pm}\) is obtained by setting spin \(i\) to \(\pm1\).
The nonnegative heat-bath operator and its symmetric form are
\begin{equation}\label{eq:finite-form}
 H_n=\sum_{i=1}^n(\sigma_i-t_i)D_i,\qquad
 \cE(f,k)=\mu\left[\sum_{i=1}^n v_iD_ifD_ik\right],
 \qquad P_t^J=e^{-tH_n}.
\end{equation}
Indeed the summand of \(H_n\) is \(f-\mu[f\mid\sigma_{-i}]\).
Conditional two-point covariance gives the displayed form. Thus
\eqref{eq:finite-form} uses the rate-one clock at each site specified
in the theorem. We write \(\cE(f)=\cE(f,f)\) and, for \(z\ge0\),
\(\cI(z)=\cE(z,\log z)\), with its extended nonnegative value at zeros.

Until the disorder comparison, \(J\) is Gaussian. Put
\(W_{ij}=J_{ij}/\sqrt n\) off the diagonal and add independent
\(W_{ii}\sim N(0,2/n)\). The diagonal contributes a constant to
\(\sigma^{\mathsf T}W\sigma/2\), so neither \(\mu\) nor \(H_n\) changes.
After independent sign choices for the eigenvectors,
\[
 W=U\diag(\lambda_1,\ldots,\lambda_n)U^{\mathsf T},\qquad
 \lambda_1\ge\cdots\ge\lambda_n,
\]
has \(U\) Haar orthogonal conditionally on the eigenvalues. Denote its
columns by \(u_a\), and put
\begin{equation}\label{eq:finite-edge}
 d_a=2-\lambda_a,\qquad g_{a,n}=T_n d_a,\qquad
 x_a(\sigma)=n^{-1/3}\langle u_a,\sigma\rangle .
\end{equation}
We use the same \(x_a\) for the coordinate functions on \(\R^\N\).
For two spins, orthogonality gives the exact identity
\begin{equation}\label{eq:replica-identity}
 n^{-2/3}\sigma\cdot\sigma'
       =\sum_{a=1}^n x_a(\sigma)x_a(\sigma').
\end{equation}
In particular, if
\(w_s=d(\nu_nP_{sT_n}^J)/d\mu\), then
\begin{equation}\label{eq:finite-spectral-traces}
 A_{n,J}(t)=\sum_{a=1}^n\mu[x_aP_{tT_n}^Jx_a],
 \qquad
 B_{n,J}(s,t)=\sum_{a=1}^n\mu[w_sx_aP_{tT_n}^Jx_a].
\end{equation}

\subsection{Conditioned measure and relaxed energy}

The following statements are the parts of the stationary manuscript
that we use directly. They also specify the meaning of the conditional
measure in the introduction. A smooth compact cylinder is a smooth
function of finitely many coordinates, compactly supported in those
coordinates.

\begin{proposition}[Stationary inputs]\label{prop:stationary-inputs}
The following hold for the edge law and the Gaussian finite models.
\begin{enumerate}
\item Let \(b+g_1>0\), \(\ell_a=(g_a+b)^{-1}\). Almost surely
      \(\sum_a\ell_a=\infty\) and \(\sum_a\ell_a^2<\infty\).
      If \(Z_a\) are independent \(N(0,\ell_a)\), the densities \(q\)
      of \(\sum_a(Z_a^2-\ell_a)\) and \(q_m\) of
      \(\sum_{a>m}(Z_a^2-\ell_a)\) are continuous, bounded, and strictly
      positive. The first \(m\) coordinates of \(\Pi_g\) have density
      \begin{equation}\label{eq:edge-cylinder-density}
       \frac{q_m\!\left(D_g(b)-\sum_{a\le m}(x_a^2-\ell_a)\right)}
            {q(D_g(b))}
       \prod_{a\le m}
          \frac{e^{-x_a^2/(2\ell_a)}}{\sqrt{2\pi\ell_a}}.
      \end{equation}
      These consistent densities are measurable in \(g\) and
      independent of \(b\). For each fixed \(0\le j<\infty\),
      \(\Pi_g[|x_a|^j]\le C_{j,g,b}\ell_a^{j/2}\), with a finite
      constant uniform in \(a\).
      For sufficiently large fixed \(m\), the tail beyond \(m\)
      has a density bounded by a finite constant relative to
      \(\bigotimes_{a>m}N(0,\ell_a)\). Conditional on that tail, the
      first \(m\) coordinates have the Gaussian angular law on the
      sphere of squared radius
      \[
       D_g(b)+\sum_{a\le m}\ell_a-\sum_{a>m}(x_a^2-\ell_a).
      \]
      In particular the renormalized constraint holds \(\Pi_g\)-almost
      surely. These are the statements of
      \cite[Proposition~2.3]{FUA}.
\item With \(b_n=1+(-g_{1,n})_+\) and \(b=1+(-g_1)_+\), the spherical
      and cube edge marginals converge jointly with the edge points to
      \(\Pi_g\) on the represented subsequences described below.
      Cube integrals converge for every continuous cylinder of
      polynomial growth. Spherical coordinates obey the uniform bounds
      \[
       \mu^{\mathrm{sp}}[|x_a|^j]\le C_j(g_{a,n}+b_n)^{-j/2},
       \qquad
       \lim_{m\to\infty}\sup_n\sum_{a>m}(g_{a,n}+b_n)^{-2}=0
      \]
      after discarding finitely many \(n\) on such a representation.
      Here \(\mu^{\mathrm{sp}}\) is surface measure on
      \(\{|\sigma|^2=n\}\), tilted by
      \(\exp(\sigma^{\mathsf T}W\sigma/2)\).
      The sources are \cite[Proposition~2.3 and Corollary~3.3]{FUA}.
\item The Euclidean gradient on smooth compact cylinders is closable
      from \(L^2(\Pi_g)\) to \(L^2(\Pi_g;\ell^2)\). The associated
      closed form has kernel equal to the constants. Every coordinate
      belongs to its domain and satisfies
      \(\nabla x_a=e_a\), hence \(\cE_g(x_a)=1\).
      This is \cite[Proposition~5.1]{FUA}. The usual contraction rule
      for gradients passes to the closure, so the form is Dirichlet
      and \(S_t^g\) is a conservative symmetric Markov semigroup.
\item On a quenched represented subsequence, say that \(f_n\) converges
      weakly to \(F\in L^2(\Pi_g)\) if its \(L^2(\mu)\) norms are bounded
      and
      \(\mu[f_n\phi(x)]\to\Pi_g[F\phi]\) for every smooth compact
      cylinder \(\phi\); call the convergence strong when in addition
      \(\mu[f_n^2]\to\Pi_g[F^2]\). Suppose, for fixed \(C,D<\infty\),
      that \(\|f_n\|_\infty\le Cn^D\), that \(T_n\cE(f_n)\) is bounded,
      and that \(f_n\) converges weakly to \(F\).
      Then \(F\) belongs to the closed form domain and
      \begin{equation}\label{eq:relaxed-liminf}
       \liminf_n T_n\cE(f_n)\ge c_*\cE_g(F).
      \end{equation}
      Conversely, for every smooth compact cylinder \(F\), there are
      functions \(f_{n,j}\), polynomially bounded for each fixed \(j\),
      such that \(f_{n,j}\) converges strongly to \(F\) as \(n\to\infty\)
      for each \(j\), and
      \begin{equation}\label{eq:relaxed-recovery}
       \lim_{j\to\infty}\lim_{n\to\infty}T_n\cE(f_{n,j})
          =c_*\cE_g(F).
      \end{equation}
      The source asserts a common represented event for countably many
      chosen sources and tests. Recovery is used at each fixed \(j\),
      followed by the displayed \(j\)-limit
      \cite[Theorem~6.8]{FUA}.
\end{enumerate}
\end{proposition}

In Item 1, the tail density is
\[
 \frac{q_m^{\mathrm{head}}\!\left(
     D_g(b)-\sum_{a>m}(x_a^2-\ell_a)\right)}{q(D_g(b))},
\]
where \(q_m^{\mathrm{head}}\) is the density of the centered head square
sum. A sufficiently large head makes this density bounded. This concrete
description will be used for estimates on infinitely many coordinates.
For a candidate pair
\((F,G)\in L^2(\Pi_g)\times L^2(\Pi_g;\ell^2)\), the angular weak
identity underlying Item 3 is
\begin{equation}\label{eq:angular-weak}
 \int \psi\,(Ax)\cdot G\,\dd\Pi_g
 =-\int F\left\{(Ax)\cdot\nabla\psi
       -\psi\sum_a(g_a+b)x_a(Ax)_a\right\}\dd\Pi_g ,
\end{equation}
for every skew-symmetric matrix \(A\) with finitely many nonzero entries
and every smooth compact cylinder \(\psi\) whose coordinate set contains
every index appearing in \(A\). These identities hold exactly when
\(F\) belongs to the closed form domain and \(G=\nabla F\)
\cite[Proposition~5.1]{FUA}. The term with \(b\) cancels by skew-symmetry.

The source also proves that \(g\mapsto A_g\), defined in
\eqref{eq:edge-semigroup}, is a Borel
\(C_{\mathrm{loc}}((0,\infty))\)-valued map and identifies \(A_g\) as
the stationary limit for both entry laws
\cite[Proposition~9.2]{FUA}. We will obtain the joint limit with
\(B_{n,J}\) from the same finite coordinates. In particular,
\eqref{eq:relaxed-liminf} is used with its polynomial bound; the
strong resolvent convergence needed here is derived in
Section~\ref{sec:gaussian-limits}.

\subsection{High-probability estimates and subsequences}

We make one convention for random constants. An estimate holds with
\emph{tight constants} if, for every \(\chi>0\), there are deterministic
constants in that estimate and events \(\mathcal G_{n,\chi}\) such that
\[
 \liminf_{n\to\infty}\Prob(\mathcal G_{n,\chi})\ge1-\chi
\]
and the estimate holds on \(\mathcal G_{n,\chi}\) simultaneously for
every scale and test function in its stated class. The constants may
depend on \(\chi\) and on fixed parameters specified before the
estimate, but not on \(n\), the scale, or the test. A tight lower
threshold for the dimension parameter \(k_p=np^3\) has the same meaning with a deterministic
threshold \(K_\chi\) on each event. When an estimate is required above
a vanishing time threshold, this means that for each \(\chi\) a
deterministic threshold \(\varepsilon_{n,\chi}\downarrow0\) can be
chosen on those events.

We use the represented-subsequence convention of
\cite[Remark~2.2]{FUA}. The edge points, the countably many cylinder
integrals and error quantities in use, and bounds that are tight in
the preceding sense are included in a joint law before applying the
subsequence and Skorokhod representation theorems. A further diagonal
subsequence handles an iterated limit such as
\(\lim_{m\to\infty}\limsup_n\). On the resulting probability space,
the relevant bounds are finite and the chosen coordinates converge
almost surely. Conditional laws of the remaining finite data recover
the original finite models. This construction is used only to show
that every subsequence has a further subsequence with the identified
limit, which implies the asserted convergence in law.

For a disorder-dependent test, we use the separate class-uniform
quenched convention of \cite[Remark~2.2]{FUA}: each invoked inequality
holds on one event for every test in its stated norm, energy, and
polynomial-growth class. The source obtains this uniformity by bounding
the error by the stated norms times a common nonnegative random error,
which permits measurable choices of the test. The common represented
event is then chosen for the countable collection of sources and tests
used in this proof. The countable core encodes limits and density classes
after the uniform bounds have been established.
